% ===== §10 Generative Asymmetry: Power Difference as a Condition of the Good =====
\section{Generative Asymmetry: When Difference of Power Generates}
\label{sec:asymmetry}

\noindent
The previous section was the paper's negative summit: it asked how the freezing of power is to be opposed, and answered that a generative relation must perpetually withhold itself from total capture. But a paper that asked only how power is prevented from doing harm would have told half the story, and the worse-lit half. For if power is the by-product of difference, and difference is the source of all generativity, then power cannot be only a threat to the good---it must, under the right conditions, be among the good's own conditions. This section is the positive summit, and it asks the question the diagnostic and defensive chapters left unasked: \emph{when, and how, does an asymmetry of power generate rather than dominate?} The answer completes the paper's account of power and supplies the criterion the praxis will need to tell cultivation from harm.

\subsection{Power and generativity are co-original}
\label{sub:asymmetry-coorigin}

The whole argument to this point has tacitly placed power and generativity in opposition: power as what damages, captures, freezes; generativity as what power threatens. The inside was generative, the outside power-laden; \emph{power-over} was the enemy, \emph{power-to} the friend. This opposition, useful as a first orientation, cannot be the last word, and the reason was given at the negative summit. Power and generativity spring from the same root. Difference generates value \emph{and} difference generates power; they are not opposed terms but co-original products of the one source. It follows that an asymmetry of power is not, as such, the antithesis of generativity. It is ambivalent---capable of freezing into domination, capable also of serving as the very condition under which certain goods are generated at all.

The point is sharpened by observing that the goods in question \emph{cannot} be generated between equals. There is no teaching between two who know the same; no nurture between two equally capable of self-care; no healing between two equally whole; no initiation between two equally versed. These are generative relations---among the most generative there are---and every one of them \emph{requires} an asymmetry of power as its precondition. To demand their symmetry is not to perfect them but to abolish them. The asymmetry is not a regrettable feature to be minimised on the way to an ideal equality; it is constitutive of the good the relation generates.

\begin{claim}[The co-originality of power and generativity]
Power and generativity are co-original products of difference, not opposed terms. An asymmetry of power is therefore ambivalent: it may freeze into domination, or it may be the constitutive condition of a good that cannot be generated between equals. Teaching, nurture, healing, and initiation are all generative relations that \emph{require} asymmetry as their precondition; to demand their symmetry is not to perfect them but to abolish them. The question is never whether asymmetry is present---in a generative relation it always is---but in which direction it runs.
\end{claim}

\subsection{The structure of good asymmetry: self-diminishing, non-appropriative}
\label{sub:asymmetry-structure}

What, then, distinguishes a generative asymmetry from a dominating one? Not its magnitude, and not its presence, but its \emph{structure in time and its relation to appropriation}. Four paradigms display the structure, and they are not chosen at random: each is a relation in which an asymmetry of power is undeniable and in which, when the relation is good, the asymmetry is exercised in a manner that tends toward its own diminution and refuses to appropriate the one it acts upon.

In \emph{teaching}, the asymmetry of knowledge is real, and the good teacher exercises it so as to dissolve it: the aim of the teaching is a student who no longer needs the teacher, an asymmetry deployed toward its own cancellation. In \emph{nurture}---the rearing of a child---the asymmetry is total at the outset and absolute in kind, and yet the whole art of good nurture is the cultivation of a being who will, in time, not require it: a power exercised for the sake of producing a subject who outgrows that power. In \emph{healing} and the tending of the vulnerable, the asymmetry of one whole and one stricken is undeniable, and the good of it lies precisely in a manner of bearing the asymmetry that does \emph{not} convert the other's weakness into the carer's standing power---that receives the vulnerability without accumulating advantage from it, and works toward the other's restored subjecthood rather than their continued dependence. In \emph{initiation and guidance}---one party more versed in some domain leading another into it---the good is the gradual mutual enrichment in which the led becomes, in their turn, one who can lead.

The common structure is unmistakable, and the series has a name for it from its earliest pages. The good asymmetry is exercised in the manner of \emph{xuande}---the ``mysterious virtue'' of the \emph{Daodejing}: to generate without possessing, to act without exacting, to lead without ruling.\footnote{\zh{生而不有，為而不恃，長而不宰，是謂玄德。} ``To give life without possessing, to act without exacting, to lead without ruling: this is called the mysterious virtue.'' See \citet{laozi1963}. The four paradigms are, severally, instances of \zh{長而不宰}---to lead or foster without ruling.} Teaching, nurture, healing, guidance: each is, when good, a \emph{zhang er bu zai}---a fostering that does not rule, a power that leads without subjugating and whose telos is the flourishing, and ultimately the empowered independence, of the one it acts upon.

\begin{claim}[The structure of generative asymmetry]
A generative asymmetry is distinguished from a dominating one not by its magnitude but by its structure: it is \emph{self-diminishing}---exercised toward its own eventual reduction rather than its perpetuation---and \emph{non-appropriative}---it does not convert the other's lesser position into standing advantage for itself. This is the power-form of \emph{xuande}, the fostering that does not rule (\zh{長而不宰}). Teaching toward the student's independence, nurture toward the child's outgrowing of it, healing toward the other's restored subjecthood, guidance toward the led one's own capacity to lead: in each, the asymmetry serves the empowerment of the one subordinate to it, and tends thereby toward its own diminution.
\end{claim}

\subsection{The criterion: power-to versus power-over}
\label{sub:asymmetry-criterion}

The structure yields a criterion, and the criterion is the affirmative counterpart of the value-theory the series has used throughout. A dominating asymmetry is \emph{power-over}: it exercises difference so as to maintain or widen the other's dependence, drawing standing advantage from the other's lesser position---which is to say, in the language of the series, treating the other as \emph{fuel}, a resource consumed to sustain the dominant party's position. A generative asymmetry is \emph{power-to}: it exercises difference so as to empower the other, such that value accrues to the one acted upon and the cycle runs toward positive holonomy. The criterion is therefore continuous with the holonomic criterion of the good cycle, and with its ethical core: \emph{does the exercise of asymmetry return value to the one subordinate to it, or does it siphon their value to the one in power?}

\begin{claim}[Power-to versus power-over]
The criterion distinguishing generative from dominating asymmetry is the direction of value's flow. \emph{Power-over} exercises difference to maintain dependence and draws advantage from the other's lesser position, reducing the other to fuel; \emph{power-to} exercises difference to empower the other, returning value to the one acted upon and running the cycle toward positive holonomy. This is the power-form of the series' standing criterion---whether anyone is reduced to mere fuel---and of the generative-justice principle that value should return to those who generate it \citep{eglash2016}. Asymmetry is just not when it is absent, which is impossible, but when its value flows toward the subordinate rather than away from them.
\end{claim}

\subsection{Trust as the mediating variable}
\label{sub:asymmetry-trust}

What decides which way a given asymmetry runs? Here the paper's third title-term returns to do its decisive work, and the two summits are joined. The variable that mediates between generative and dominating asymmetry is \emph{trust}.

Consider the four paradigms again. In each, the asymmetry generates the good \emph{only because it is mediated by trust}: the student entrusts themselves to the teacher in the confidence that the knowledge-power will not be turned to keep them dependent; the child's utter exposure to the parent's total power is borne by---and only safely borne by---a trust that the power will be exercised for the child's outgrowing of it; the patient yields their vulnerability to the healer on trust that the weakness will not be converted into the healer's leverage. Subtract the trust, and the identical asymmetry inverts: the teacher who is not trusted to relinquish power becomes the master who hoards it; the unentrusted parental power becomes control; the carer who exploits rather than receives the vulnerability becomes the abuser. \emph{Same asymmetry, opposite outcome---and trust is the hinge.}

\begin{claim}[Trust as the hinge of asymmetry]
Trust is the variable that decides whether an asymmetry of power becomes generative or frozen. Asymmetry suffused with trust empowers: the one in the lesser position can entrust themselves to the greater power in the confidence that it will be exercised toward their empowerment and not their subjection, and on that confidence the good---the teaching, the nurture, the healing, the guidance---is generated. The same asymmetry voided of trust dominates: the lesser position becomes exposure to be exploited, and the power-to collapses into power-over. This is why trust stands in the title beside power; and it is the precise point at which the present paper welds to its predecessor's account of how trust is built \citep{paperxiii}, for the production of trust there described is, here, revealed as the condition under which difference of power generates rather than dominates.
\end{claim}

\subsection{The completion of the account of power}
\label{sub:asymmetry-completion}

The two summits now stand together, and the account of power is complete. The negative summit (\xref{sec:balance}) showed that power must be prevented from freezing---that a generative relation withholds itself from total capture, perpetually guarding the lines of flight against the second-order mechanisms that would seize them. The positive summit shows that power, unfrozen and mediated by trust, is not merely to be tolerated but is among the conditions of the good---that the very asymmetry the negative summit guards against freezing is, when it runs toward the other's empowerment, the generative source of teaching, nurture, healing, and guidance. The two are not in tension; they are the two faces of a single thesis about difference. Difference generates power; the power can freeze, and must be kept from freezing; the same power, kept liquid and turned toward the other, generates the good. The whole normative content of the paper is the maintenance of that double condition: keep the asymmetry from freezing, and keep it turned toward empowerment---which is to say, keep it suffused with trust and incompletely captured.

This completes the answer to the second of the paper's three questions---how love is not destroyed by power---by showing that the question was incompletely posed. Power does not only threaten love; rightly structured and rightly trusted, difference of power is part of how love generates a world at all. There remains the third question, the one neither the diagnosis of harm nor the account of asymmetry has yet touched: how a bond, having learned to generate the good and to guard it against freezing, might come to be \emph{free}---free not in the abstract but within history, against the inherited weight of the material conditions that laid down its asymmetries in the first place. To that question, and to the dialectical-materialist account it requires, the paper now turns.
